% =====================================================================
%  Abstract della tesi.
%
%  Il testo e' quello del Sommario del Lavoro di Tesi
%  (file_latex_sommario/sezioni/*.tex), adattato all'inserimento qui:
%
%    - titolazioni non numerate: prima del Capitolo 1 una \section
%      numerata uscirebbe come "0.1" e finirebbe nell'indice;
%    - equazioni e tabella in numerazione araba semplice, ripristinata
%      subito dopo, cosi' il Capitolo 1 riparte da 1.1;
%    - etichette prefissate "abs-", per non collidere con i capitoli;
%    - tabella con [H] (pacchetto float): un float differito verrebbe
%      composto dopo il ripristino dei contatori e si rinumererebbe;
%    - nessuna voce in elenco figure/tabelle;
%    - bibliografia propria rimossa: le citazioni confluiscono
%      nella bibliografia unica in fondo alla tesi (le chiavi sono gia'
%      tutte in Bibliografia.bib).
%
%  ATTENZIONE: questo file e' una COPIA del sommario, non un include.
%  Overleaf riceve solo file_latex/, quindi non si poteva puntare alla
%  cartella del sommario. Il testo va allineato nel sommario e nelle due tesi.
% =====================================================================

% --- pagina divisoria ------------------------------------------------
\clearpage
\thispagestyle{empty}
\phantomsection
\addcontentsline{toc}{chapter}{Abstract}
\vspace*{\fill}
\begin{center}
    {\Huge\bfseries Abstract}
\end{center}
\vspace*{\fill}
\clearpage

% --- numerazione locale ----------------------------------------------
\renewcommand{\theequation}{\arabic{equation}}
\renewcommand{\thetable}{\arabic{table}}
\let\abstractACL\addcontentsline
\renewcommand{\addcontentsline}[3]{}

\begingroup
\setlength{\parindent}{0pt}
\setlength{\parskip}{4pt plus 0.5pt minus 0.5pt}

\section*{Introduzione}

L'algoritmo di Shor~\cite{shor1994} permette a un computer quantistico di
scomporre un numero nei suoi fattori, come $15=3\times5$.
In teoria, per numeri molto grandi richiede meno risorse di calcolo del
miglior metodo classico noto~\cite{lenstra1993gnfs}.

Questa possibilità mette a rischio sistemi crittografici come RSA,
la cui sicurezza si basa sulla difficoltà di trovare quei fattori~\cite{rsa1978}.
Nei dispositivi quantistici attuali, però, gli errori fisici possono
compromettere il calcolo: qui sono indicati come \emph{rumore}~\cite{preskill2018}.

Le dimostrazioni di Shor restano limitate a piccoli numeri~\cite{skosana2021}.
Questa tesi usa simulazioni per studiare come scegliere meglio i risultati
delle misure, correggere gli errori e semplificare il circuito.
Un modello di apprendimento automatico valuta i risultati delle misure
per scegliere quali analizzare. Il confronto senza il modello,
chiamato \emph{ablazione}, mostra se questo offre un vantaggio.
Esperimenti separati studiano come proteggere l'informazione quantistica
dagli errori. Una rete neurale cerca di riconoscerli e scegliere la correzione.
Si semplifica anche la trasformata di Fourier quantistica (QFT), che
aiuta a ricavare i fattori. Eliminare alcune operazioni riduce il
circuito, ma può rendere il risultato meno preciso.

Le prove iniziali hanno rivelato un errore nel circuito simulato.
Il circuito è stato corretto e verificato prima di ripetere gli esperimenti.
La verifica controlla anche l'operazione realizzata, perché i soli picchi corretti non bastano a dimostrare la correttezza del circuito.

\section*{Materiali e Metodi}

\subsection*{Come si ricavano i fattori}

Si sceglie un numero $a$ senza fattori in comune con $N$.
Il circuito cerca il \emph{periodo}: dopo quanti passi si ripetono i resti
delle potenze di $a$ divise per $N$.
Una procedura classica usa il periodo per ricavare e verificare i
fattori~\cite{nielsenChuang2000}.

Una singola esecuzione con misura finale è detta \emph{shot}.
Nel caso principale, $N=15$ e $a=7$, otto qubit, le unità di informazione
quantistica, permettono di misurare 256 valori.
Senza errori, si ottengono quattro valori: $0$, $64$, $128$ e $192$.
Tre consentono di trovare i fattori, quindi il successo ideale per
misura è $3/4=75\%$.
Il rumore può disperdere questi picchi e rendere meno affidabile la scelta del solo esito più frequente.

Con esiti tutti equiprobabili, la procedura classica ne accetta comunque
63 su 256. Questo successo casuale, detto \emph{pavimento uniforme}, vale circa il $24{,}61\%$.
Per questo, il successo della verifica classica non basta da solo a dimostrare che il circuito abbia conservato informazione quantistica utile.

\subsection*{Rumore e mitigazione}

Il rumore viene applicato alle operazioni elementari del circuito,
dette \emph{porte}, tenendo conto della loro durata.

Il modello include depolarizzazione, che rende lo stato più casuale,
perdita di energia e perdita delle proprietà quantistiche.
Comprende anche errori di lettura delle misure.
Si confrontano UC1 (riferimento) e UC2, con depolarizzazione cinque volte più intensa
e perdita di energia e delle proprietà quantistiche più rapida.

La \emph{zero-noise extrapolation}~\cite{temme2017} ripete il calcolo
con rumore più intenso per stimare il risultato senza rumore.
Le due varianti considerate usano due o tre livelli di rumore.

\subsection*{Classificazione e confronto statistico}

Ogni ciclo raccoglie $1\,024$ misure in un \emph{istogramma},
che conta quante volte compare ciascun esito.
Il Metodo~1 controlla l'esito più frequente; TOP-4 i quattro più frequenti.
Il Metodo~2 usa TOP-4 solo quando il modello prevede il successo sul primo esito.
Confrontarlo con TOP-4 senza filtro separa il contributo del classificatore da quello della ricerca di più candidati.

Si generano $2\,000$ istogrammi per scenario, variando il rumore entro il $\pm50\%$ del riferimento.
Sono divisi fra addestramento ($60\%$), selezione ($20\%$) e verifica finale ($20\%$).
Si confrontano random forest (alberi decisionali), macchine a vettori
di supporto (SVM) e reti neurali.

F1 valuta il riconoscimento degli istogrammi utili e la correttezza delle
previsioni di successo. L'AUC misura quanto bene si separano dati utili e inutili.
Valori vicini a uno sono migliori; si sceglie la SVM con F1 più alto.

Si contano i cicli necessari per trovare i fattori, fino a 50.
Ai fallimenti si assegna il valore 51.
I metodi ricevono gli stessi istogrammi.
I test statistici usano $p_{\mathrm{Holm}}$, corretto per i confronti multipli:
sotto $0{,}05$ le differenze sono considerate significative.

\subsection*{Correzione d'errore e decoder}

La correzione d'errore (QEC) protegge un \emph{qubit logico}, l'unità di
informazione da conservare, con più qubit fisici.
I controlli segnalano gli errori senza leggere l'informazione protetta.
L'insieme di questi segnali è chiamato \emph{sindrome}.
Il \emph{decoder} interpreta questi segnali e sceglie la correzione.
$p$ indica la probabilità di errore fisico e $p_L$ quella che
l'informazione resti sbagliata dopo la correzione.

Si studiano i codici a ripetizione, Steane~\cite{steane1996} e surface code~\cite{fowler2012}.
Steane usa sette qubit per uno logico; il surface code li dispone su un reticolo bidimensionale.
La distanza $d$ è il minimo numero di qubit coinvolti in un errore
che altera l'informazione senza essere rilevato.

Nel surface code, gli errori colpiscono anche preparazione e misura.
Il \emph{matching}, metodo di riferimento, abbina i segnali dei controlli
per individuare gli errori più probabili.
Si stima la soglia $p_{\mathrm{th}}$ confrontando distanze $d=3,5,7,9$.
Sotto questa soglia di rumore, aumentare la distanza riduce l'errore dopo la correzione.

La rete è valutata sia da sola sia come aiuto al matching.
Nel decoder ibrido, la rete interviene solo quando prevede che la decisione del matching debba essere invertita.
I decoder ricevono gli stessi segnali; le scelte sono fissate prima della verifica finale.

\subsection*{Sensibilità di Shor e QFT approssimata}

Una prova separata inserisce errori dopo ogni porta di Shor con probabilità $p_g$.
Non usa codici correttori: $p_g$ è distinto dall'errore dopo la correzione $p_L$.

Su $N=21$ si eliminano piccole rotazioni della QFT, sia durante i calcoli
aritmetici sia nella trasformata finale.

Test separati stimano la fase, un angolo associato a un'operazione quantistica.
Si usano sei, otto e dieci qubit, con tre fasi per dimensione.
È corretta la misura più vicina alla fase cercata, alla precisione disponibile.
L'approssimazione è scelta su dati separati dalla verifica finale.
Il rumore deriva da FakeSherbrooke, modello simulato di un dispositivo IBM.
Le varianti usano la stessa disposizione dei qubit.

\subsection*{Strumenti e validazione}

Si usano Qiskit/Aer~\cite{qiskit2023} per i circuiti, Stim~\cite{stim2021}
e PyMatching~\cite{pymatching2021} per la QEC, scikit-learn~\cite{scikit_learn} per i modelli appresi.

I circuiti aritmetici di Beauregard~\cite{beauregard2002} per $N=21,35$
sono verificati controllando risultati e qubit ausiliari, usati nei calcoli intermedi.
Per $N=21$, la distanza di variazione totale fra risultati simulati ideali
e teorici è $1{,}38\times10^{-9}$, indicando un accordo quasi esatto.

Codice, dati e versioni sono disponibili nel
\href{https://github.com/claudio-dragotta/shor-nisq-qec-thesis}{repository della tesi}.

\section*{Risultati}

\textbf{Classificazione.} Il circuito verificato ha profondità 412, cioè 412 strati
di porte in sequenza. Contiene 224 porte CX, che invertono un qubit in base al valore di un altro.
Nei due scenari UC1/UC2, le SVM raggiungono F1 $=0{,}919/0{,}923$ e AUC $=0{,}965/0{,}958$.
Su 30 repliche per scenario, tutte le strategie fattorizzano entro il limite previsto.

La sola ricerca TOP-4 porta i cicli medi a $1{,}00$, contro $1{,}37/1{,}50$ del Metodo~1.
Il confronto è significativo ($p_{\mathrm{Holm}}=0{,}0024/0{,}0015$).

Il Metodo~2 richiede $1{,}50/1{,}37$ cicli e non migliora significativamente il Metodo~1
($p_{\mathrm{Holm}}=0{,}7314/0{,}0874$). Richiede più cicli di TOP-4
($p_{\mathrm{Holm}}=0{,}0033/0{,}0096$).
Buone metriche di classificazione non garantiscono quindi una riduzione del costo necessario a trovare i fattori.

La zero-noise extrapolation non migliora il riferimento.
La variante su tre livelli triplica il costo di campionamento.

\textbf{Correzione d'errore.} Il codice a ripetizione a $d=3$ riproduce i valori previsti dalla teoria.
Per Steane, l'errore dopo la correzione cresce circa con il quadrato dell'errore fisico (esponente stimato $1{,}94$).
La pseudo-soglia di Steane è $p\simeq0{,}08$: qui gli errori dopo la
correzione eguagliano quelli del qubit non protetto.
Riguarda un codice a distanza fissata.

Per il surface code, le soglie stimate sono $0{,}86\%$ ($Z$) e $0{,}81\%$ ($X$).
$Z$ e $X$ indicano due modi di preparare e misurare il qubit.

La rete usata da sola non supera il matching nelle quattordici configurazioni iniziali ($d=3,5$).

Con crosstalk (interferenza fra qubit), la rete affiancata al matching riduce $p_L$ del $18$--$21\%$ a $d=3$ e dell'$1$--$3\%$ a $d=5$.
A $d=7$ non emerge un miglioramento significativo.
In questi test, alcune correlazioni non sono rappresentate dal grafo del matching e lasciano un margine per la rete.

\textbf{Sensibilità di Shor.} Il successo senza rumore è $74{,}89\%$, con intervallo al 95\%
$[74{,}59\%;\ 75{,}18\%]$, coerente con il $75\%$ teorico.

Il successo si stabilizza intorno a $p_g=0{,}1$; a $p_g=0{,}5$ vale $24{,}59\%$.
È vicino al pavimento uniforme del $24{,}61\%$, dovuto agli esiti accettati anche senza informazione utile.

\textbf{QFT approssimata.} Su $N=21$ eliminare le piccole rotazioni durante
i calcoli aritmetici riduce le porte a due qubit da $49\,661$ a $23\,178$.
Il successo senza rumore scende però dal $46{,}67\%$ al $40{,}13\%$.
Nelle simulazioni rumorose eseguite non emerge un vantaggio misurabile.

Un modello semplificato stima un guadagno massimo di circa tre punti percentuali.
Richiede errori delle porte a due qubit circa duecento volte inferiori alla calibrazione considerata.

Nella stima di fase, tutti i nove confronti favoriscono la QFT troncata.
In sei confronti, l'intervallo al 95\% del miglioramento è interamente positivo.
Un intervallo interamente positivo distingue il guadagno da zero nel singolo confronto.
Il maggior aumento è di 12,6 punti percentuali su dieci qubit
(intervallo al 95\%: 9,8--15,3 punti).
La scelta migliore conserva le rotazioni con gli angoli più grandi,
limitandoli a due o tre valori nelle dimensioni esaminate.

\section*{Discussione e Conclusioni}

L'ablazione mostra che la riduzione dei cicli dipende dalla ricerca TOP-4.
Il classificatore può scartare dati utili.
TOP-4 è efficace nel caso studiato, ma non è una scelta ottima per ogni istanza.

Nella correzione d'errore, la rete aiuta il matching quando sono presenti
interferenze fra qubit. Il beneficio diminuisce con la distanza del codice.

Tutti gli esperimenti sono simulazioni. Lo studio del classificatore usa
solo $N=15$, con rumore uniforme e invariato nel tempo.
Per il classificatore, gli errori di operazioni diverse sono indipendenti; non si considerano
interferenze fra qubit né operazioni aggiuntive richieste dalle loro connessioni.

La correzione considera soltanto alcune operazioni e alcuni tipi di errore:
le stime possono quindi essere favorevoli~\cite{plaquette2026}.
Non viene simulato un circuito completo di Shor protetto dai codici.

Il vantaggio della QFT approssimata nella stima di fase non si estende
automaticamente a Shor. Le prove rumorose su $N=21$ confrontano varianti
della QFT; $N=35$ è verificato senza rumore.
La stima di fase usa una disposizione dei qubit per dimensione,
tre fasi e un solo insieme di dati di calibrazione.

Occorre collegare gli errori dopo la correzione alle operazioni di Shor
ed estendere i modelli di rumore.
Per i decoder, si potranno usare reti che considerino disposizione dei
qubit ed evoluzione dei controlli~\cite{alphaqubit2024}.
La QFT approssimata andrà combinata con metodi aritmetici che riducono
i qubit necessari~\cite{chevignard2025}, confrontando porte, qubit e profondità.

\par
\endgroup

% --- ripristino della numerazione della tesi -------------------------
\let\addcontentsline\abstractACL
\renewcommand{\theequation}{\thechapter.\arabic{equation}}
\renewcommand{\thetable}{\thechapter.\arabic{table}}
\setcounter{equation}{0}
\setcounter{table}{0}
\clearpage
